\section{Final-separator descent and normalized local angular gauges}
\label{sec:r13-final-separator-raw-descent}

The completed marked curve will be kept as the \emph{old} coefficient
chart when the last source-separating functions are added.  This order
matters.  Contracting new sources only to the tangent plane of the
pre-completion chart does not control mixed relations involving its
new Picard coordinates.

\subsection{A linear descent statement for the entire labelled star}

Let $T$ and $N$ be finite-dimensional real vector spaces.  Let
$v_p\in N\setminus\{0\}$, $p\in M$, be a finite labelled family
such that
\[
 \ker\bigl(\mathbf R^M\longrightarrow N,\ (a_p)\longmapsto
                    \sum_pa_pv_p\bigr)=\mathbf Rw,
       \qquad w_p>0.
 \tag{SD1}
\]
Write $W_p=T\oplus\mathbf Rv_p$.  On a connected convex open
base $B\subset T$, let $b_p$ be continuous finite piecewise smooth
covectors on the respective half-products, with common tangential
strain.  Their traces are denoted $\beta_p=b_p|_T$ and their
primitive normal traces $c_p=b_p(v_p)$.  Fix $p_0\in M$ and put
\[
 \beta_0=\beta_{p_0},\qquad
 K_p=\beta_p-\beta_0,\qquad z_w=\sum_pw_pc_p.
 \tag{SD2}
\]
By CS6 every $K_p$ is affine Killing on $B$.

Consider a second finite labelled incident system
$W'_\alpha\supset T'$ in an ambient space $W'$, and a linear
map
\[
 \pi:W'\longrightarrow T\oplus N,
                    \qquad \pi|_{T'}:T'\xrightarrow{\sim}T.
 \tag{SD3}
\]
Partition its labels into sets $I_p$, $p\in M$, and $I_0$ such
that
\[
 \pi(W'_\alpha)\subset W_p\quad(\alpha\in I_p),
   \qquad \pi(W'_\alpha)\subset T\quad(\alpha\in I_0).
 \tag{SD4}
\]
No independence among the final spaces or between these sets is
assumed.  Labels with zero old normal image remain in $I_0$.
Choose the pulled-back traces
\[
 b'_\alpha=\pi^*b_p\quad(\alpha\in I_p),\qquad
 b'_\alpha=\pi^*\beta_0\quad(\alpha\in I_0).
 \tag{SD5}
\]
The second expression is defined on the relevant ordinary image
contained in $T$, rather than by restricting an arbitrary current
to $T$.

\begin{theorem}[Every mixed raw Hessian descends to the old flux]
\label{thm:r13-sd-mixed-raw-hessian}
If $\operatorname{Hess}_Tz_w=0$, then every raw defect for the
\emph{whole} final labelled system has zero full Hessian on
$B'=(\pi|_{T'})^{-1}B$.  This includes simultaneous relations
mixing different old marks with arbitrary silent labels.
The conclusion also holds for any other continuous finite-PS
branch primitives with the same strains as SD5.
\end{theorem}

\begin{proof}
Take $u'=(u'_\alpha)$ with $\sum_\alpha u'_\alpha=0$.
For $\alpha\in I_p$ write uniquely
\[
 \pi u'_\alpha=t_\alpha+a_\alpha v_p,
 \qquad A_p=\sum_{\alpha\in I_p}a_\alpha,
 \tag{SD6}
\]
and put $t_\alpha=\pi u'_\alpha$ for $\alpha\in I_0$.
Applying $\pi$ to the relation gives
$\sum_\alpha t_\alpha=0$ and $\sum_pA_pv_p=0$.
Thus $A=cw$ for a real constant $c$.  The final raw defect is
exactly
\[
 \begin{split}
 \sum_\alpha b'_\alpha(u'_\alpha)
   &=\sum_pA_pc_p+
       \sum_p\sum_{\alpha\in I_p}K_p(t_\alpha)\\
   &=c z_w+
       \sum_p\sum_{\alpha\in I_p}K_p(t_\alpha).
 \end{split}
 \tag{SD7}
\]
The terms involving the common trace $\beta_0$ cancel because
$\sum t_\alpha=0$.  The remaining evaluations of the fixed
vectors $t_\alpha$ in affine Killing fields are affine.  This
proves the asserted Hessian vanishing, with every mixed relation
retained.  The difference of two primitives with the same strain
on a connected convex half-product is affine Killing by CS6 in
all its variables.  Its boundary evaluation on $u'_\alpha$ is
affine, proving the last assertion.
\end{proof}

The statement requires no decomposition of the final raw kernel
into separate source kernels.  In particular it is stronger than
testing each balanced source independently.  Conversely it
requires SD4 relative to the \emph{completed} old columns: making
silent labels tangent only after forgetting those columns is
insufficient for SD6.

\subsection{An explicit normalized gauge on one whole positive face}

Assume now $T=T_0\oplus E$, with $\dim E=2$, and let
$B=B_0\times\operatorname{relint}(\gamma)$, where $B_0\subset T_0$
is connected and convex and $\gamma\subset E$ is a convex
two-dimensional cone.  All functions extend continuously to
$B_0\times\{0\}$ and satisfy the central normalization
\[
        b_p(t,0,0)|_{T_0}=b_{T_0}(t),\qquad
        b_p(t,0,0)|_{E\oplus\mathbf Rv_p}=0,
 \tag{LG1}
\]
with one common covector $b_{T_0}(t)$; it need not be zero.
Their regularity is the closed-cell regularity of CS.  Assume
$\operatorname{Hess}_Tz_w=0$ on the \emph{whole} $B$.

\begin{theorem}[A local normalized angular gauge]
\label{thm:r13-lg-normalized-angular-gauge}
There are constant alternating forms on the three-dimensional
normal spaces $E\oplus\mathbf Rv_p$, extended by zero on
$T_0$, whose associated linear Killing corrections make every
raw incident defect zero along $B$.  These corrections preserve
the strains and LG1.  After this correction, the branch
covectors admit a common ambient covector trace on $B$.
Pulling the corrected primitives through SD3--SD5 gives the same
conclusion for the whole final system.
\end{theorem}

\begin{proof}
By CS6, $K_p$ is an affine Killing covector on $T_0\oplus E$.
Its value is zero on $B_0\times\{0\}$ by the common central
value in LG1.  Write
$K_p(y)=a_p+\Omega_py$ with $\Omega_p$ alternating.  Vanishing
on the nonempty open tangent set implies $a_p=0$ and
$\Omega_pT_0=0$.  Alternation also removes all mixed
$T_0$--$E$ entries.  Thus
\[
              K_p(t,x)=\Omega_p^E x,
                   \qquad \Omega_p^E\in\bigwedge^2E^*.
 \tag{LG2}
\]
Add the Killing covector $-\Omega_p^E x$ on branch $p$,
extended independently of its primitive transverse variable
$s$.  This makes all tangential traces equal to $\beta_0$.
It leaves every $c_p$ and hence $z_w$ unchanged.

The affine function $z_w$ vanishes at $(t,0)$ for every
$t\in B_0$, again by LG1.  Consequently
\[
                              z_w(t,x)=\ell(x)
                         \quad\text{for some }\ell\in E^*.
 \tag{LG3}
\]
On one chosen branch $p_0$, add the further Killing covector
\[
       \kappa_{p_0}(x,s)
           ={s\over w_{p_0}}\,\ell
                  -{\ell(x)\over w_{p_0}}\,ds.
 \tag{LG4}
\]
Here $s$ is the primitive parameter with tangent vector $v_{p_0}$.
Its coefficient matrix is alternating, so its strain is zero.
At $s=0$ its tangential trace is zero, while its contribution to
the weighted normal flux is $-\ell(x)$.  All corrected branches
therefore have the same tangential trace and flux zero.  Every
raw relation vanishes by CS15 and SD1.

The restriction map from the ambient dual of
$T+\operatorname{span}\{v_p\}$ to the finite direct sum of the
$W_p^*$ has image exactly the annihilator of its raw sum kernel.
The corrected traces belong to that image.  A fixed linear
right inverse on this image supplies one ambient covector trace,
continuous and finite-PS in the base variables.  For the final
system use its pullback through $\pi$; each silent branch takes
the common tangent restriction.  All corrections vanish at
$(x,s)=0$, so LG1 is preserved.
\end{proof}

The gauge is local to the one positive face.  Different faces
may require different constant alternating forms on a shared
three-cone.  The existence of a single set of cone gauges valid
at all faces is still the separate JF8 equation.  Local gauges
on every nonzero angular neighborhood do not prove that equation.

\subsection{Finite projection of a supplied actual curve block}

The following use of projection is a calculation on one original
analytic source.  It is not a proper pushforward statement for
an arbitrary current.  Let a fixed adapted tuple $a$ have
radial rank two on an actual curve block supplied by FC and CC.
For calculation, enlarge it by the two CC normal parameters and
a finite generating family for $k(C)$.  These are actual units
on the same ordinary neighborhood.  CC identifies the whole
source of this calculating tuple as the Gauss product, with
the marked calibration $\delta_pm_p$ at every radius.

\begin{lemma}[Calibrated finite projection without a new global completion]
\label{lem:r13-sd-finite-calculating-projection}
The complete top-dimensional source contribution of $a$ on this
ordinary block is the coordinate projection of that entire
calibrated calculating source.  A branch of projected rank less
than three contributes zero to every top-degree form pulled
back from $a$.  A full-rank projected branch has the factor
\[
             i_T\,{m_p^{a}\over m_p},
             \qquad
        \delta_pm_p\,i_T\,{m_p^{a}\over m_p}
                    =\delta_pi_Tm_p^{a},
 \tag{SD8}
\]
where $i_T$ is the saturated tangent lattice index and
$m_p^{a}$ the stretching in the saturated old normal quotient.
All coincident branch images are summed only after their labels
and these factors have been retained.
\end{lemma}

\begin{proof}
Retention of $a$ means that every full-rank $a$ source is a
full-rank source of the calculating tuple.  CC identifies all
the latter sources on the specified ordinary block.  The
sourcewise coordinate-refinement formula retained in GP and
CC18 therefore computes the full $a$ contribution.  A
top-degree differential form pulled back through a map of rank
less than three is literally zero on that branch; no current
restriction is being used.  On a full-rank branch, the exact
sequence of tangent and normal lattices gives SD8, precisely as
in CC18.  The ordinary domain and the same analytic source are
unchanged by this calculation.  The calculating coordinates
need not be added to the global tuple on every other source.
\end{proof}

\subsection{Actual detector after one final separator}

\begin{lemma}[Finite normal Newton preparation after the last separator]
\label{lem:r13-sd-finite-normal-newton}
Fix an SP completed curve block with its two rational normal
parameters $z_1,z_2$.  A further fixed finite list of rational
functions in $F$ is adapted after a finite rational conical
subdivision of the positive two-face and regular toroidal
modification.  More precisely, remove the finite one-dimensional
angular tie rays of its normal Newton diagrams.  Each remaining
two-dimensional sector has one leading normal monomial for
each numerator and denominator, with coefficient in
$k(C)^\times$.  The regular two-cones of a finite refinement
are attached to smooth proper curves $C'$ with
$k(C')=k(C)$, hence $C'\simeq C$.  On each such curve the
new functions are adapted, and the old normalized marked orders
are unchanged under this identification and the old coordinate
projection.  Any new mark outside the old marked support has
zero old completed normal column.
\end{lemma}

\begin{proof}
Write each rational function as $A/B$ with nonzero
$A,B\in\mathcal O_{V,\eta_C}$.  Complete this regular
two-dimensional local ring and choose a coefficient field.
It is $k(C)[[z_1,z_2]]$.  The support of each nonzero
series has finitely many minimal elements for the coordinatewise
partial order on $\mathbf N^2$: induction on the first
coordinate, retaining only successive strict decreases of the
second, proves this elementary finiteness.  Every exponent
dominates a minimal one.  Therefore every strictly positive
weight takes its minimum on that finite set.  The finitely
many equal-weight rays divide the positive quadrant into
sectors with a unique leading exponent for every $A$ and $B$.
Taking the union over the finite function list gives one finite
rational fan.

Perform the regular toroidal subdivision used in FC.  On each
regular two-cone the pullback of a series is its selected
normal monomial times a unit: every other exponent differs
by a nonnegative vector in the dual cone, and at least one
of its two normal exponents is positive.  The same assertion
holds in the local ring, by faithful flatness of its completion.
Taking the quotient gives the adapted expression for $A/B$.
Over the generic point of $C$, the new two-cone is a split
toric fixed point.  Its residue field is precisely $k(C)$.
Its curve closure is thus birational to $C$; the proper smooth
curve identification follows.  Monomial normal parameters
can be chosen from the original rational $z_1,z_2$ on these
generic charts.  The old normalized residue functions are
then the same functions of $k(C)$.  If other trivializations
are used, SH gives the equivalent common shear, and the old
normal coordinate projection is unchanged.

In particular all old normalized residue functions have order
zero at a point outside their old divisor supports.  New
functions may make such a point visible, but its old completed
normal column remains zero.  Their own finite supports add
only finitely many marks.  The tie rays and all subdivision
boundary rays are lower-dimensional faces; no conclusion of
the two-dimensional-sector statement is assigned to them.
\end{proof}

For an actual analytic unit of norm one at the central point,
first use its rational residue representative.  Their ratio
has residue one, so a strict ordinary inequality gives equal
norms and normalized initial forms on a common central
neighborhood.  This is the finite replacement step of FC.
The lemma therefore applies to the final FA and protection
coordinates without assuming in advance that they were adapted
to the old curve.

We formulate the geometric combination on the supplied
three-dimensional central model of FC, with any supplied
calibrated old tangent factor.  Start with finitely many
specified original two-faces.  Prepare their proper smooth
marked curves by FC.  Apply SP simultaneously and retain the
resulting completed tuple $j$.  Each specified block now has
its old labelled columns $v_p$, weights $w_p=\delta_pm_p$, and
kernel SD1.  All original coefficient coordinates are retained.

\emph{Only after this completion}, apply FA to the whole compact
full source $S_j$, and retain its finite separating coordinates
in a tuple $h$.  Any further finite protection coordinates are
retained too.  Choose a common finite source/target subdivision
for this final tuple.  On the positive two-face pieces considered
below its $j$ projection has rank two and lies in the specified
old two-face.  Work in one sufficiently short common collar;
every face in the conclusion is a whole connected positive
prepared piece of this collar, not an isolated compact box.

\begin{theorem}[Complete positive-face detector with the last separator]
\label{thm:r13-sd-actual-final-separator-detector}
Assume the original coefficient tensor $Z$ has continuous
finite-PS branch primitives as in CS, with common tangential
strain, and that $D[Z]=0$ against every original compact smooth
complementary CLD test.  Use the protected complete calibrated
source presentation of GP, the finite face preparation and
rank alternative of FC, and the all-radius comparison of CC.
Then, over each of the above positive rank-two pieces:
\begin{enumerate}
\item the full distributional Hessian of the old completed
      flux $z_w$ vanishes, including all angular, old tangent,
      and mixed entries;
\item every raw Hessian for the \emph{entire final incident
      system}, including mixed core/satellite relations and
      all zero-old-image labels, vanishes;
\item if the fixed branch system extends to the central face
      with LG1, there is the normalized local angular gauge
      of LG2--LG4 and a continuous finite-PS ambient covector
      lift on a neighborhood of that positive face.
\end{enumerate}
The finite calculating enlargements in the preceding lemma are
used only to evaluate source contributions.  They do not
require a further global Picard completion or a repeated
separator/completion process.
\end{theorem}

\begin{proof}
Fix a compact box in a positive prepared piece.  FC is applied
with \emph{old tuple $j$}.  Because $h$ separates the complete
old full source, every final full-rank source over this box is
either the same actual old core source or has its entire old
$j$ ordinary image in the tangent plane.  This is the strict
ordinary containment of FC4, including commensurable radii.
It is stronger than a dimension statement at a source point.
The full protected source sum is finite, and its source germs
are complete balanced cycles by GP and local balancing.

At the old core, CC and a finite adapted calculation give the
complete marked normal cycle.  Any final incident branch projects
either to one of its old labelled rays or to zero in its old
normal quotient.  AH supplies an explicit finite angular horn
for the finished adapted list supplied by
Lemma~\ref{lem:r13-sd-finite-normal-newton}; a further common finite source
subdivision keeps each old coefficient branch specified.
If a normal parameterization needs extra curve generators,
Lemma~\ref{lem:r13-sd-finite-calculating-projection} computes its
contribution on the same ordinary block.  The full lattice
projection formula gives, in compatible primitive coordinates,
\[
         i_T\sum_{\alpha\mapsto p}w'_\alpha\lambda_\alpha=w_p,
       \qquad \pi(v'_\alpha)=\lambda_\alpha v_p,
                      \quad\lambda_\alpha>0.
 \tag{SD9}
\]
Silent labels have $\pi(v'_\alpha)=0$.  The tangent factor
$i_T>0$ is fixed on the prepared piece and is not silently
absorbed into the old weights.

Use the pullback of the old branch primitive on every branch
projecting to an old mark.  On each entire satellite source,
FC4 makes the represented tensor the literal pullback of the
common old tangential strain; use $\pi^*\beta_0$ there.
Choose the quotient splitting through the retained old tangent
projection, so every silent normal vector has zero old image.
Its normal flux is zero.  Its tangential Green terms cancel on
its complete balanced cycle.  The unequal tangential traces of
the main branches cancel by CS10--CS11.  Thus CS, SD9, and the
complete original current equation give
\[
                  \operatorname{Hess}_T(i_T^{-1}z_w)=0
 \tag{SD10}
\]
on the box.  Every symmetric tangent dual is paired with this
same full source calculation, so no marginal spanning argument
replaces the paired-symbol test.  The factor $i_T^{-1}$ is
nonzero and constant.  In SD10 the tangent variables have first
been transported through $\pi|_{T'}$ and all symmetric duals
in the final tangent variables are varied.  Equivalently, if
one uses the pullbacks of the old tangent test forms, their
tangent integration Jacobian supplies the factor $i_T$ and
SD10 reads directly as $\operatorname{Hess}_Tz_w=0$ in old
coordinates.  Both conventions give the same full Hessian
vanishing, proving the first assertion.

For every final raw relation apply SD6--SD7.  The silent spaces
are tangent \emph{for $j$}, which is the tuple whose old marked
kernel is SD1.  Hence the theorem on mixed raw Hessians applies
to the entire system without assuming that its kernel splits
as a sum of separate source kernels.  It also makes the
conclusion independent of the chosen same-strain primitives.

The boxes cover the whole positive prepared piece in the same
fixed finite coefficient system.  Distributional Hessian
vanishing is local, so it holds on that whole piece; its affine
expressions agree on overlaps.  With the central continuity and
normalization assumed in the third assertion, LG applies.  Its
common ambient trace extends over each incident half-product by
a fixed linear complement to its tangent restriction.  These
extensions agree on the face.  A finite source subdivision gives
a continuous finite-PS covector on the local union, with the
original branch strains.  Pullback through the final coordinate
projection gives the stated final-system lift.
\end{proof}

\begin{corollary}[Coverage of a complete original two-face inventory]
\label{cor:r13-sd-complete-original-two-face-inventory}
Suppose the original finite coefficient tuple $f$ has first been
protected by its complete-source faithful tuple as in FC, and
the supplied original curve list contains \emph{every} labelled
old full-source two-face under consideration.  Retain this
first faithful tuple through SP and through the last separator
for $j$.  Then every final positive two-face whose $f$ projection
has rank two and is contained in one of those original geometric
two-faces has a local continuous finite-PS covector lift, under
the regularity and central normalization of the theorem.
\end{corollary}

\begin{proof}
The intermediate $j$ face is either a completed original core
piece, or its actual source is not one of the listed old full
sources.  In the first case apply the theorem.  In the second
case FC, now with old tuple $f$ and the first faithful tuple
retained, places the entire ordinary $f$ image in the original
two-plane.  The represented tensor is therefore the literal
pullback of its common tangential strain.  Pick one tangential
primitive and pull it back; this is already a local primitive
on that entire ordinary source.  After the last separator for
$j$, all other final sources over that intermediate face have
their ordinary $j$ images in its tangent plane.  Their $f$
images are in the same original two-plane.  The same pulled-back
primitive therefore works on their complete union, including
every mixed relation.

All labels are accounted for because the original full-source
inventory is complete and FC prepares every remaining finite
source piece.  No assertion about only a selected old source is
substituted for this hypothesis.  A final face projecting into
the interior of one old top cone is simpler: the one primitive
on that cone pulls back throughout its local coefficient image.
\end{proof}

\paragraph{Exact scope.}
The source hypothesis is now a proved combination of the stated
FC, SP, CC, AH, FA and GP inputs for the supplied model and its
finite positive rank-two face list.  It is not asserted for a
general ordinary nonsaturated analytic germ without those model
inputs.  New one-dimensional angular walls, faces whose old
projection has rank at most one, and the common central face
retain their own raw equations.  In particular the theorem
does not prove JF8 for the entire three-fan.  It gives no
ordinary-to-GK upgrade, no universal higher-degree relation
descent, and no general A1 statement.

\paragraph{Proof provenance.}
SD1 and the old weights are SP4 and CC15--CC18; complete original
tests and the calibrated source sum are GP and its retained
sourcewise coordinate-refinement formula.  FC is used with the
\emph{post-SP} tuple as its old tuple.  AH controls the fixed
adapted finite lists on their original ordinary horns.  The
new algebraic steps are SD6--SD7 and the explicit local gauge
LG2--LG4.  The Stokes identity used throughout is the separately
reviewed CS theorem, which includes all internal coefficient
walls and all mixed differential terms.
